\chapter{Methodology}
\label{ch:methodology}

\section{Experimental Design}
The benchmarking methodology relies heavily on controlled synthetic workloads to isolate configuration effects. A Python-based workload generator synthesizes transaction requests using a seeded Poisson arrival process. The system then executes the workload through the sequencer, executor, and submitter pipeline, recording telemetry at each defined boundary.

\section{Data Collection and Metrics}
Telemetry data is collected as raw JSON/CSV metrics in the \texttt{benchmark-suite/metrics/} directory. The Data Tools analytics pipeline subsequently processes this data using Python aggregation scripts to compute the final metrics.

\begin{table}[H]
\centering
\caption{Key Benchmarking Metrics}
\label{tab:metrics}
\small
\begin{tabularx}{\textwidth}{l X c X}
\toprule
\textbf{Metric} & \textbf{Definition} & \textbf{Unit} & \textbf{Why it matters} \\
\midrule
Throughput & Rate of successful transaction inclusion & TPS & Indicates capacity limits of the system \\
Latency & End-to-end time from transaction submission to L1 finalization & ms & Affects user experience and finality \\
Gas Cost & L1 execution cost per transaction or per batch & Wei & Determines the economic viability of the rollup \\
\bottomrule
\end{tabularx}
\end{table}

\section{Configuration Sweeps}
Experiments are conducted by methodically varying:
\begin{itemize}
    \item \textbf{Batch Size:} Evaluating small vs. large batches to observe the cost-amortization vs. queueing-delay effect.
    \item \textbf{Data Availability Mode:} Comparing standard \texttt{calldata} versus EIP-4844 blobs.
    \item \textbf{Mock Proof Delay:} Utilizing parameterized proof delays to isolate the impact of proof generation time from standard sequencing and submission overhead.
\end{itemize}
